% =============================================================================
% Module 16: Solutions to Exercises
% =============================================================================
% This file is conditionally included from main.tex via \ifsolutions
% =============================================================================

\section{Solutions to Exercises}
\label{sec:oi_solutions}

\noindent
Full worked solutions are also available in the Mathematica script
\solnlink{16_Odd_Images_Central/problems\_16.wl}{problems\_16.wl},
which verifies each answer symbolically and numerically.

% ----- Exercise 16.1 -----
\subsection*{Exercise~\ref{ex:oi_gradient_hessian}: Fermat Gradient and Hessian}

\textbf{(a)} The geometric term of $\tau$ has gradient
$\vectheta-\vecbeta$ and the potential term contributes
$-\nabla\psiL=-\vecalpha$, so
\[
  \nabla_{\vectheta}\tau = (\vectheta-\vecbeta)-\vecalpha ,
\]
and $\nabla_{\vectheta}\tau=0$ is exactly $\vecbeta=\vectheta-\vecalpha$.
Solving $\nabla\tau=0$ for the source returns
$\vecbeta=\vectheta-\nabla\psiL$ identically.

\textbf{(b)} Differentiating once more, the geometric term gives the
identity and the potential term $-\partial_a\partial_b\psiL$:
\[
  \nabla\nabla\tau = \mathbb{I}-\nabla\nabla\psiL
  = \begin{pmatrix}1-\psiL_{,11} & -\psiL_{,12}\\ -\psiL_{,12} & 1-\psiL_{,22}\end{pmatrix}
  = \pderiv{\vecbeta}{\vectheta} = \Amat ,
\]
and $\operatorname{tr}\Amat = 2-\nabla^2\psiL = 2(1-\kappab)$.

\textbf{(c)} In the eigenbasis $\Amat=\operatorname{diag}(\lambda_-,\lambda_+)$
with $\lambda_\pm=1-\kappab\pm|\gammaL|$, so
$\det\Amat=\lambda_+\lambda_-$ and $\operatorname{tr}\Amat=\lambda_++\lambda_-$.
Two real numbers are both positive iff their product and sum are
positive, both negative iff the product is positive and the sum
negative, and of opposite sign iff the product is negative.  Hence:
minimum $\Leftrightarrow$ $\det\Amat>0$, $\operatorname{tr}\Amat>0$;
maximum $\Leftrightarrow$ $\det\Amat>0$, $\operatorname{tr}\Amat<0$;
saddle $\Leftrightarrow$ $\det\Amat<0$, as in Table~\ref{tab:morse_types}.
\emph{[Attested: \solnlink{16_Odd_Images_Central/problems\_16.wl}{problems\_16.wl}, checks 1a--1h.]}

% ----- Exercise 16.2 -----
\subsection*{Exercise~\ref{ex:oi_five_images}: Five Images of a Cored Elliptical Lens}

A source inside the tangential caustic has four bright images: two
minima outside the tangential critical curve and two saddles, which
contribute $2-2=0$ to the index sum.  The cored lens satisfies
(S1)--(S2), so by eq.~\eqref{eq:oi_index_sum} the indices must add up to
$1$: there is at least one more positive-parity image.  Since the source
is also inside the radial caustic, it is the maximum of $\tau$ supported
by the core, so $N=5$ with $n_{\min}=2$, $n_{\rm sad}=2$, $n_{\max}=1$:
\begin{center}
\begin{tabular}{clcc}
\toprule
\# & Fermat type & signs of $(\lambda_-,\lambda_+)$ & $\operatorname{sgn}\mu$\\
\midrule
1, 2 & minimum (Type~I)   & $(+,+)$ & $+1$\\
3, 4 & saddle (Type~II)   & $(-,+)$ & $-1$\\
5    & maximum (Type~III) & $(-,-)$ & $+1$\\
\bottomrule
\end{tabular}
\end{center}
and $\sum\operatorname{sgn}\mu_i=2+1-2=1$, as required.

\textbf{Numerical check.}  For
$\psiL=\sqrt{0.25^2+0.7\,\theta_1^2+1.3\,\theta_2^2}$ and
$\vecbeta=(0.05,0.03)$, solving $\vecbeta=\vectheta-\nabla\psiL$ with
\texttt{NSolve} gives exactly five real images.  Classifying each by
the eigenvalues of $\Amat=\mathbb{I}-\nabla\nabla\psiL$ gives two minima
($\mu=2.19$ and $2.35$), two saddles ($\mu=-1.51$ and $-1.21$), and one
maximum at $|\vectheta|=0.029$ with $\mu=0.136$.  The maximum is both
the image closest to the centre and the faintest, and the parities sum
to $+1$.
\emph{[Attested: \solnlink{16_Odd_Images_Central/problems\_16.wl}{problems\_16.wl}, checks 2a--2h.]}

% ----- Exercise 16.3 -----
\subsection*{Exercise~\ref{ex:oi_nis_central}: The NIS Central Image}

\textbf{(a)} From eq.~\eqref{eq:oi_nis},
$\alpha(\theta)=\thetaE\theta/(\sqrt{\theta^2+\theta_c^2}+\theta_c)
=\tfrac{\thetaE}{2\theta_c}\,\theta+\mathcal{O}(\theta^3)$, so
$\alpha/\theta\to\kappab_0$ and $\dd\alpha/\dd\theta\to\kappab_0$ as
$\theta\to0$, with $\kappab_0=\thetaE/(2\theta_c)$.  By
eq.~\eqref{eq:oi_lambdas}, $\lambda_t(0)=\lambda_r(0)=1-\kappab_0$.  Both
are negative iff $\kappab_0>1$, i.e.\ iff $0<\theta_c<\thetaE/2$; then
the centre (an image for $\ysrc=0$) is a maximum of $\tau$.

\textbf{(b)}
\[
  \mu_c = \frac{1}{\lambda_t(0)\lambda_r(0)} = \frac{1}{(1-\thetaE/2\theta_c)^2}
        = \frac{4\theta_c^2}{(\thetaE-2\theta_c)^2}
        = \frac{4\theta_c^2}{\thetaE^2} + \mathcal{O}(\theta_c^3).
\]
For $\thetaE=1$: $\mu_c=1/16=0.0625$ at $\theta_c=0.1$ and
$\mu_c=1/2401=4.16\times10^{-4}$ at $\theta_c=0.01$.

\textbf{(c)} Keeping the linear term of the lens map,
$\ysrc\simeq(1-\kappab_0)\,\theta$, so the central image lies at
$\theta\simeq\ysrc/(1-\kappab_0)=-\ysrc/(\kappab_0-1)$: on the side of
the lens opposite the source.  For $\theta_c=0.1\,\thetaE$
($\kappab_0=5$) and $\ysrc=0.01\,\thetaE$ this is
$\theta\simeq-0.0025\,\thetaE$; the exact root of
eq.~\eqref{eq:oi_nis_cubic} agrees to better than $2\times10^{-6}\,\thetaE$.

\textbf{(d)} As $\theta_c\to0$, $\alpha(\theta)\to\thetaE$ for every
$\theta>0$: the SIS.  Then $\mu_c\to0$, and the SIS magnification along
the far side, $\mu=1/(1-\thetaE/|\theta|)$ (with $\lambda_r=1$), tends to
zero as $\theta\to0$.  The would-be central image has been pushed into
the conical tip of the SIS with zero flux: the SIS shows two images,
not three.
\emph{[Attested: \solnlink{16_Odd_Images_Central/problems\_16.wl}{problems\_16.wl}, checks 3a--3j.]}

% ----- Exercise 16.4 -----
\subsection*{Exercise~\ref{ex:oi_puncture}: Point Masses Change the Parity of the Count}

\textbf{(a)} On a circle of radius $\epsilon$ around a point mass,
$\nabla\tau = \vectheta-\vecbeta-\vecalpha_{\rm smooth}-\thetaBH^2\hat{\vectheta}/\epsilon$.
As $\epsilon\to0$ the last term dominates every other (they stay
finite), so $\nabla\tau$ points radially inward all the way round: it
winds once, index $+1$, independent of $\vecbeta$.  Numerically, the
winding is $+1$ for $\ysrc=0.1,\,0.5,\,2,\,5\;\thetaBH$.

\textbf{(b)} Draw a large circle (winding $+1$) and excise a small disc
around each point mass.  The winding around the large circle equals the
sum of the image indices plus the indices of the $g$ punctures:
$n_+-n_-+g=1$, which is eq.~\eqref{eq:oi_index_pm}.  Then
$N=n_++n_-=2n_-+1-g$.  Adding one point mass changes $1-g$ by one, so
it flips the parity of $N$; and $N+g=2n_-+1$ is always odd.  Counting
each puncture as a phantom image restores Burke's odd total.

\textbf{(c)} For $\theta_c=0.2\,\thetaE$ and
$\thetaBH^2=0.3\,\theta_{\bullet,\rm crit}^2$: at $\ysrc=0.1\,\thetaE$ the
NIS alone has 3 images and the NIS with the black hole has 2; at
$\ysrc=0.2\,\thetaE$ the counts are 3 and 4.  Both changes are odd
($-1$ and $+1$), and with $g=1$, $N+g=3$ and $5$ are odd.  In the first
case the central maximum has merged with the black-hole saddle (the
source is inside the inner caustic); in the second it survives and the
black-hole saddle has been added.
\emph{[Attested: \solnlink{16_Odd_Images_Central/problems\_16.wl}{problems\_16.wl}, checks 4a--4c.]}

% ----- Exercise 16.5 -----
\subsection*{Exercise~\ref{ex:oi_sis_tip}: The Tip of the SIS}

\textbf{(a)} On a circle of radius $\epsilon\to0$,
$\nabla\tau\to-\vecbeta-\thetaE\hat{\vectheta}$.  As $\hat{\vectheta}$
goes once around the unit circle, the point $-\thetaE\hat{\vectheta}$
traces a circle of radius $\thetaE$ about the origin, and the field is
this circle shifted by $-\vecbeta$.  It winds once around zero iff the
shifted circle encloses the origin, i.e.\ iff $|\vecbeta|<\thetaE$;
otherwise its winding is zero.  Numerically: winding $1$ for
$\ysrc=0.2,\,0.5,\,0.9\,\thetaE$ and $0$ for $\ysrc=1.1,\,1.5,\,3\,\thetaE$.

\textbf{(b)} For $|\vecbeta|<\thetaE$ the SIS has two images,
$\theta_+=\ysrc+\thetaE$ with $\mu_+>0$ and $\theta_-=\ysrc-\thetaE$ with
$\mu_-<0$: image indices sum to $0$, plus tip index $1$, total $1$.
For $|\vecbeta|>\thetaE$ there is one image, $\mu>0$: sum $1$, plus tip
index $0$, total $1$.  The winding at infinity is always $1$, as it must
be.  The circle $|\vecbeta|=\thetaE$ in the source plane --- the
\emph{cut} of the SIS --- plays the role of a caustic: crossing it
changes the image count by one, not two, because the second ``image''
is the singular tip, which carries index but no flux.
\emph{[Attested: \solnlink{16_Odd_Images_Central/problems\_16.wl}{problems\_16.wl}, checks 5a--5b.]}

% ----- Exercise 16.6 -----
\subsection*{Exercise~\ref{ex:oi_critical_mass}: The Critical Black-Hole Mass}

\textbf{(a)} Expanding $(\thetaE\theta_c^2/2)^{2/3}=2^{-2/3}\thetaE^{2/3}\theta_c^{4/3}$,
\[
  \theta_{\bullet,\rm crit}^2 = \thetaE\theta_c - 3\cdot2^{-2/3}\,\thetaE^{2/3}\theta_c^{4/3} + \theta_c^2 ,
\]
exactly.  Dividing by $\thetaE\theta_c$, the second and third terms
scale as $(\theta_c/\thetaE)^{1/3}$ and $\theta_c/\thetaE$, so the ratio
tends to $1$ as $\theta_c/\thetaE\to0$.  The approach is slow: the
leading correction is $-3\cdot2^{-2/3}(\theta_c/\thetaE)^{1/3}$, which
is $-0.70$ at $\theta_c=0.05\,\thetaE$.

\textbf{(b)} For $\theta_c=0.05\,\thetaE$:
$\theta_{\bullet,\rm crit}^2=0.01769\,\thetaE^2$, so
$\theta_{\bullet,\rm crit}=0.133\,\thetaE$.

\textbf{(c)} The projected mass inside $\theta$ scales as
$\theta\,\alpha(\theta)$ for any axisymmetric lens, and as
$\thetaBH^2$ for a point mass, with the same proportionality constant.
For $\theta_c=0.1\,\thetaE$, $\alpha(\thetaE)=(\sqrt{1.01}-0.1)\,\thetaE=0.905\,\thetaE$,
so the critical mass fraction is
$\mBH=0.02228/0.905=0.0246$: a black hole with about $2.5\%$ of the
mass inside the Einstein radius erases every central image of this lens.
\emph{[Attested: \solnlink{16_Odd_Images_Central/problems\_16.wl}{problems\_16.wl}, checks 6a--6e.]}

% ----- Exercise 16.7 -----
\subsection*{Exercise~\ref{ex:oi_powerlaw}: A Shallow Cusp}

\textbf{(a)} With $\gp=1.5$ and $\thetaE=1$, $\alpha=\operatorname{sgn}(\theta)|\theta|^{1/2}$.
On the far side, $\theta=-x=-w^2$ ($w>0$) gives $\ysrc=w-w^2$, with roots
\[
  w = \frac{1\mp\sqrt{1-4\ysrc}}{2},
\]
both positive iff $0<\ysrc<1/4$ (the maximum of $w-w^2$ is $1/4$ at
$w=1/2$, which is the radial caustic).  On the near side,
$\theta=w^2$ gives $\ysrc=w^2-w$, which has exactly one positive root.
So a source at $\ysrc<\thetaE/4$ has three images: an odd number, as
the continuity of $\alpha$ guarantees.

\textbf{(b)} The central image is the smaller far-side root,
$w_c=(1-\sqrt{1-4\ysrc})/2$.  With $\lambda_t=1-1/w$ and
$\lambda_r=1-1/(2w)$,
\[
  \mu_c = \frac{2w_c^2}{(w_c-1)(2w_c-1)} .
\]
At $\ysrc=0.1$: $w_c=0.1127$ and $\mu_c=0.0370$, while
eq.~\eqref{eq:oi_mu_powerlaw} gives $2\ysrc^2=0.02$.  The leading-order
law underestimates by a factor $1.85$, because the corrections are of
relative order $w_c\simeq\ysrc$.

\textbf{(c)} For $\gp=1.8$, $\alpha=|\theta|^{0.2}$ on each side, and
the far-side equation $x^{0.2}-x=0.1$ is solved numerically.  The
central image has $\mu_c=5.007\times10^{-8}$, which agrees with
the leading-order value $5\times10^{-8}$ to $0.2\%$: here the correction
is of relative order $\ysrc^{4}$.  The image is $18.3$ magnitudes
fainter than the unlensed source.  Moving the slope from $1.5$ to
$1.8$ costs almost six orders of magnitude in central-image flux at
fixed source position: the steep cusp has suppressed it, as
Section~\ref{sec:oi_steep} argued.
\emph{[Attested: \solnlink{16_Odd_Images_Central/problems\_16.wl}{problems\_16.wl}, checks 7a--7d.]}
